\clearpage
\thispagestyle{plain}
\section*{Abstract}
\phantomsection
\addcontentsline{toc}{section}{Abstract}

Urban-mobility simulators inform high-stakes infrastructure decisions in
transportation planning, congestion management, and policy evaluation,
yet a foundational methodological gap persists: cross- simulator results
are not directly comparable because each simulator uses different input
formats, different execution conventions, and different metric
definitions. A 10 \% difference in mean travel time between two
simulators on a nominally identical scenario could reflect genuine
paradigm-level disagreement or any of several input-interpretation,
configuration-tuning, version-pinning, floating-point, threading, and
platform-binary effects. Practitioners who must choose a simulator for
an actual planning decision have no operational way to distinguish
these.

This thesis presents \textbf{SimForge}, a reproducible cross-simulator
testing framework for urban mobility simulation that closes this gap
through four contributions: (C1) a canonical, simulator-agnostic input
bundle with field-level validation; (C2) byte-deterministic adapters for
SUMO, MATSim, and DTALite that translate the canonical bundle into each
engine\textquotesingle s native input format and apply a shared
strongly-connected-component + feasibility filter; (C3) a pinned-digest
container image (\texttt{ghcr.io/phanidharakula/simforge}) published to
GitHub Container Registry via an automated GitHub Actions workflow,
providing bit-identical reproducibility across machines; and (C4) a
programmatic cross-engine fairness audit that verifies byte-identical
feasibility verdicts (Q1), byte-identical SCC networks (Q2), identical
simulated trip-count targets (Q3), and a cross-engine mean travel-time
comparison (Q4), together with a reproducibility index R = 1 − σ/μ and
95 \% confidence intervals on every reported mean.

The framework integrates three simulator paradigms (microscopic queue
via SUMO, activity-based queue via MATSim, user-equilibrium dynamic
traffic assignment via DTALite) and is benchmarked across three U.S.
metropolitan networks (Chicago, New York City, Los Angeles) at demand
tiers from 1 K to 500 K trips per scenario. The fairness audit (Q1-Q3)
PASSes at every shipped tier, empirically demonstrating that the
canonical-input + deterministic-adapter design isolates engine-paradigm
differences from input-interpretation differences.

Four empirical findings emerged from the build and contribute to the
cross-simulator benchmarking literature beyond the original framework
deliverables. First, a shared canonical-routes BFS module (Phase 14)
reduced the chicago\_200k\_car benchmark wall by approximately 20 ×
cold-vs-cold (from 141.87 h to 7.14 h on OSC Cardinal) and 228 × for
warm-cache re-runs, making the 500 K-trip tier (nyc\_500k\_car)
tractable for the first time (12 h 8 min wall vs an \textasciitilde600 h
pre-Phase-14 projection). Second, the SUMO ↔ MATSim mean travel-time
ratio is regime-dependent: small-tier convergence (within ±5 \% at 50 K
trips) sharply reverses at saturation, reaching a 96.3 \% gap at 500 K
trips on NYC\textquotesingle s bridge-and-corridor topology, where
SUMO\textquotesingle s insertion-refusal model and
MATSim\textquotesingle s queue-hold model engage divergent
paradigm-level behaviors that the fairness contract makes attributable
to mobsim choice rather than input asymmetry. Third, at a fixed code
version the framework is bit-identical across every measured platform:
20/20 cells byte-identical between Cardinal x86\_64 host venv and
Cardinal x86\_64 container (different OS, JDK distribution, and JDK
major version); Mac arm64 ↔ Cardinal x86\_64 agreement to 4 sig figs on
mean travel time. The 0.2-3 \% per-engine shifts originally observed in
cross-platform comparisons are attributable to code-version drift
between time-separated measurements (Phase 12 ↔ Phase 14.13
routing-layer changes), not to floating-point hardware differences. The
pinned-digest container is therefore load-bearing because it freezes the
code + bundle + toolchain at a precise git SHA, not because it closes
any cross-platform numerical gap. Fourth --- sibling to the second
finding in mechanism but orthogonal in axis --- the within-engine SUMO
micro vs meso mean-TT gap \emph{narrows} at saturation (from +27 \% at 1
K trips to +7.8 \% at 200 K) while the wall-time premium grows
monotonically (1.65 × → 9 × → \textasciitilde124 ×). At saturation, both
mobsim resolutions become bound by the same insertion-refusal dynamics
that drive the cross-engine divergence; lane-level dynamics that micro
adds become second-order to the bottleneck behavior. Two paradigm-spread
phenomena (cross-engine, within-engine) thus emerge from the same
underlying mechanism, surfaced cleanly only because the fairness
contract holds.

SimForge is open-source under the Apache License 2.0 and available at
\texttt{github.com/PhanidharAkula/SimForge}. The framework, the
canonical bundles, the audit tooling, and the pinned-digest container
are intended as a foundation for future cross-simulator research that
can build on a bit-reproducible benchmark target rather than re-deriving
fairness from scratch for each comparison.

\textbf{Keywords:} urban mobility simulation, cross-simulator
benchmarking, reproducibility, SUMO, MATSim, DTALite, dynamic traffic
assignment, agent-based simulation, container reproducibility,
scientific computing.
